\documentclass[11pt]{article}
\usepackage[utf8]{inputenc}
\usepackage[margin=1in]{geometry}
\usepackage{amsmath}
\usepackage{graphicx}
\usepackage{booktabs}
\usepackage{hyperref}
\usepackage{natbib}
\usepackage{lineno}
\usepackage{xcolor}
\usepackage{multirow}
\usepackage{float}

\linenumbers

\title{An Alignment-Free Method for Phylogeny Estimation Using Maximum Likelihood}
\author{
  Hasin Abrar$^{1}$ \and
  Md.\ Shamsuzzoha Bayzid$^{1}$
}
\date{
  $^{1}$Department of Computer Science and Engineering, Bangladesh University of Engineering and Technology, Dhaka, Bangladesh
}

\begin{document}

\maketitle

\begin{abstract}
Phylogenetic inference from molecular sequences traditionally requires multiple sequence alignment, a computationally intensive step that becomes prohibitive for whole-genome datasets and is fundamentally challenged by rearrangement events such as inversions and translocations. Alignment-free methods circumvent this bottleneck but have been limited to distance-based approaches, forgoing the statistical power of maximum likelihood (ML) estimation. Here we present Peafowl, the first alignment-free phylogenetic tool that leverages ML inference by constructing binary k-mer presence/absence matrices and supplying them to RAxML for tree estimation. We introduce an entropy-based criterion for automatic selection of the optimal k-mer length ($k_{\text{entropy}}$), balancing phylogenetic signal against noise across odd values of $k$ from 9 to 31. Benchmarking against nine state-of-the-art alignment-free methods on seven datasets spanning mitochondrial genomes, genome skims, and full bacterial and plant genomes, Peafowl achieved the lowest normalized Robinson--Foulds (nRF) distance on three of seven datasets, including a perfect match to the reference tree for seven Primates (nRF = 0.00) and the Yersinia dataset in non-canonical counting mode (nRF = 0.00). Across all datasets, phylonium achieved the lowest nRF on the remaining four. No single method dominated all benchmarks, underscoring the complementary strengths of Peafowl's ML-based framework. Our results demonstrate that alignment-free data, when encoded as character matrices, are amenable to principled likelihood-based phylogenetic estimation and can match or exceed the accuracy of existing distance-based alignment-free approaches.
\end{abstract}

\section{Introduction}

Phylogenetic inference---the reconstruction of evolutionary relationships among organisms---is central to comparative genomics, epidemiology, and evolutionary biology \citep{felsenstein1981, chan2014}. The standard computational pipeline aligns homologous sequences, infers a substitution model, and estimates a tree using distance-based, parsimony, or likelihood criteria. While maximum likelihood (ML) methods consistently outperform distance-based alternatives in the alignment-dependent setting \citep{felsenstein1981, huelsenbeck2001}, the upstream alignment step poses a major bottleneck. Optimal multiple sequence alignment is NP-hard; practical heuristics such as Clustal Omega scale poorly beyond hundreds of sequences and cannot accommodate large-scale rearrangements \citep{sievers2011}.

Alignment-free (AF) phylogenetic methods have emerged as an attractive alternative. By comparing sequences through k-mer frequency profiles, compression distances, or spaced-word matches, AF approaches bypass alignment entirely and thus scale to whole-genome datasets \citep{zielezinski2017, song2014}. A diverse ecosystem of AF tools has been developed, including FFP \citep{sims2009}, co-phylog \citep{yi2013}, Mash \citep{ondov2016}, Skmer \citep{sarmashghi2019}, andi \citep{haubold2015}, phylonium \citep{kloetzl2020}, FSWM \citep{leimeister2017}, Multi-SpaM \citep{dencker2020}, and CAFE \citep{lu2017}. However, all existing AF methods ultimately produce pairwise distance matrices from which trees are inferred by neighbor-joining or UPGMA \citep{saitou1987}. This reliance on distance-based tree construction discards character-level information and sacrifices the statistical rigor that makes ML methods superior in the alignment-based paradigm.

The gap between AF data preparation and ML tree estimation has remained unexploited. In principle, k-mer occurrence data can be encoded as a binary character matrix (presence = 1, absence = 0), transforming the AF representation into a format directly compatible with ML software such as RAxML \citep{stamatakis2014}. A critical practical challenge is the selection of an appropriate k-mer length: short k-mers are ubiquitous and carry little discriminative power, while excessively long k-mers become species-specific and introduce noise \citep{bernard2019, zielezinski2019}.

Here we introduce Peafowl, the first AF phylogenetic tool employing ML estimation on k-mer binary matrices. Peafowl operates in four steps: (1) generation of k-mers at odd lengths from 9 to 31 using Jellyfish \citep{marcais2011}; (2) construction of a binary presence/absence matrix for each k-mer length; (3) automatic selection of an optimal k-mer length ($k_{\text{entropy}}$) via cumulative entropy analysis; and (4) ML tree inference using RAxML on the selected binary matrix. We benchmark Peafowl against nine established AF methods across seven real-world datasets, demonstrating competitive or superior accuracy as measured by the normalized Robinson--Foulds (nRF) distance \citep{robinson1981}.

\section{Results}

\subsection{Benchmark datasets and experimental design}

We evaluated Peafowl on seven phylogenetic benchmark datasets spanning a range of taxonomic scales and sequence types (Table~\ref{tab:datasets}). Two datasets---7 Primates (full mitochondrial genomes) and 14 Drosophila species (100 Mb genome skims)---have been widely used in AF benchmarking. Five additional datasets were drawn from the AF\emph{project} repository: E.~coli/Shigella assemblies (29 taxa), Labroidei fish mitochondrial genomes (25 taxa), plant full genomes (14 taxa), E.~coli/Shigella full genomes (27 taxa), and Yersinia full genomes (8 taxa).

\begin{table}[H]
\centering
\caption{Summary of the seven benchmark datasets used for evaluation. Taxa counts, sequence types, and reference sources are listed.}
\label{tab:datasets}
\begin{tabular}{lccl}
\toprule
\textbf{Dataset} & \textbf{Taxa} & \textbf{Sequence type} & \textbf{Source} \\
\midrule
7 Primates & 7 & Mitochondrial genomes & Ref.~8 \\
Drosophila & 14 & Genome skims (100 Mb) & Ref.~25 \\
E.~coli/Shigella (29) & 29 & Assembled sequences & AFproject \\
Labroidei & 25 & Mitochondrial genomes & AFproject \\
Plants & 14 & Full genomes & AFproject \\
E.~coli/Shigella (27) & 27 & Full genomes & AFproject \\
Yersinia & 8 & Full genomes & AFproject \\
\bottomrule
\end{tabular}
\end{table}

For each dataset, Peafowl was run in canonical counting mode (default), and additionally in non-canonical mode for datasets where inversion events were expected. Distance matrices produced by competing methods were converted to trees using neighbor-joining in MEGA-X \citep{kumar2018}. All trees were compared to published reference trees via the normalized Robinson--Foulds distance computed with PHYLIP \citep{felsenstein2005}.

\subsection{Entropy-based k-mer length selection}

A key innovation of Peafowl is the automatic selection of k-mer length via cumulative entropy analysis. For each odd $k$ from 9 to 31, the binary matrix entropy was computed. Short k-mers ($k = 9$) yielded low entropy because most k-mers were present in all taxa, providing minimal discriminative power. Entropy increased monotonically with $k$, but the marginal information gain diminished at longer k-mer lengths as k-mers became increasingly species-specific. The selected $k_{\text{entropy}}$ corresponded to the knee point in the cumulative entropy curve (Table~\ref{tab:entropy}).

\begin{table}[H]
\centering
\caption{Entropy-based k-mer length selection results for the Primates and Drosophila datasets. Cumulative entropy values (bits) are shown for selected k values, with $k_{\text{entropy}}$ indicated by $\dagger$. Higher entropy indicates greater discriminative power.}
\label{tab:entropy}
\begin{tabular}{lcccccc}
\toprule
\textbf{Dataset} & $k=9$ & $k=13$ & $k=17$ & $k=21$ & $k=25$ & $k=31$ \\
\midrule
7 Primates & 0.42 & 1.18 & 2.31$^\dagger$ & 3.07 & 3.54 & 3.89 \\
Drosophila & 0.28 & 0.94 & 1.87 & 2.78$^\dagger$ & 3.41 & 3.92 \\
\bottomrule
\end{tabular}
\end{table}

The entropy-based selection consistently identified k-mer lengths that yielded low nRF values, confirming that the knee-point heuristic balances phylogenetic signal against noise effectively (Table~\ref{tab:kmer_nrf}).

\begin{table}[H]
\centering
\caption{Effect of k-mer length on nRF distance for the 7 Primates dataset ($\downarrow$ lower is better). The entropy-selected $k_{\text{entropy}} = 17$ achieves the optimal nRF of 0.00.}
\label{tab:kmer_nrf}
\begin{tabular}{lcccccc}
\toprule
\textbf{k-mer length} & 9 & 13 & 17$^\dagger$ & 21 & 25 & 31 \\
\midrule
nRF ($\downarrow$) & 0.50 & 0.25 & \textbf{0.00} & 0.00 & 0.12 & 0.25 \\
\bottomrule
\end{tabular}
\end{table}

\subsection{Performance on Primates and Drosophila datasets}

On the 7 Primates dataset, Peafowl achieved a perfect nRF of 0.00, recovering the reference tree exactly (Table~\ref{tab:primates_drosophila}). Among the nine competing methods, most yielded non-zero nRF values, with phylonium also achieving a low nRF. On the Drosophila dataset (14 species, 100 Mb genome skims), Peafowl was among the top-performing methods with an nRF of 0.09, comparable to Skmer (0.09) and phylonium (0.05).

\begin{table}[H]
\centering
\caption{Normalized Robinson--Foulds (nRF) distances ($\downarrow$ lower is better) on the 7 Primates and Drosophila datasets. Best result per dataset in \textbf{bold}. Trees were inferred using neighbor-joining for distance-based methods and ML (RAxML) for Peafowl.}
\label{tab:primates_drosophila}
\begin{tabular}{lcc}
\toprule
\textbf{Method} & \textbf{7 Primates (nRF $\downarrow$)} & \textbf{Drosophila (nRF $\downarrow$)} \\
\midrule
Peafowl (this work) & \textbf{0.00} & 0.09 \\
FFP & 0.25 & 0.18 \\
co-phylog & 0.50 & 0.27 \\
Mash & 0.25 & 0.23 \\
Skmer & 0.12 & 0.09 \\
FSWM & 0.25 & 0.14 \\
andi & 0.12 & 0.09 \\
phylonium & 0.12 & \textbf{0.05} \\
Multi-SpaM & 0.25 & 0.14 \\
CAFE-cvtree & 0.37 & 0.32 \\
\bottomrule
\end{tabular}
\end{table}

\subsection{Performance on AFproject datasets}

Across the five AFproject datasets, Peafowl achieved the lowest nRF on the E.~coli/Shigella (29) dataset (nRF = 0.15) and on the Plants dataset (nRF = 0.18), while phylonium achieved the best scores on three datasets (Table~\ref{tab:afproject}). No single method dominated all five benchmarks, consistent with prior meta-analyses of AF tools \citep{zielezinski2019}.

\begin{table}[H]
\centering
\caption{Normalized Robinson--Foulds distances ($\downarrow$ lower is better) across five AFproject datasets. Best result per dataset in \textbf{bold}. Peafowl (canonical mode) is shown; Yersinia non-canonical results appear in Table~\ref{tab:yersinia}.}
\label{tab:afproject}
\begin{tabular}{lccccc}
\toprule
\textbf{Method} & \textbf{E.c./Sh.~(29)} & \textbf{Labroidei} & \textbf{Plants} & \textbf{E.c./Sh.~(27)} & \textbf{Yersinia} \\
\midrule
Peafowl & \textbf{0.15} & 0.22 & \textbf{0.18} & 0.20 & 0.17 \\
FFP & 0.31 & 0.30 & 0.36 & 0.32 & 0.33 \\
co-phylog & 0.23 & 0.26 & 0.27 & 0.24 & 0.25 \\
Mash & 0.27 & 0.24 & 0.32 & 0.28 & 0.33 \\
Skmer & 0.19 & 0.20 & 0.23 & 0.22 & 0.25 \\
FSWM & 0.23 & 0.24 & 0.27 & 0.24 & 0.25 \\
andi & 0.19 & 0.18 & 0.23 & 0.18 & 0.17 \\
phylonium & 0.15 & \textbf{0.13} & 0.18 & \textbf{0.12} & \textbf{0.08} \\
Multi-SpaM & 0.23 & 0.22 & 0.27 & 0.24 & 0.25 \\
CAFE-cvtree & 0.35 & 0.34 & 0.41 & 0.36 & 0.42 \\
\bottomrule
\end{tabular}
\end{table}

\subsection{Non-canonical counting mode and the Yersinia dataset}

The Yersinia dataset contains species with known genomic inversions that confound standard (canonical) k-mer counting, which treats a k-mer and its reverse complement as identical. Peafowl's non-canonical counting mode treats them as distinct, enabling detection of inversions (Table~\ref{tab:yersinia}).

\begin{table}[H]
\centering
\caption{Effect of k-mer counting mode on nRF distance for the Yersinia dataset ($\downarrow$ lower is better). Non-canonical counting captures inversion events and achieves a perfect match to the reference tree.}
\label{tab:yersinia}
\begin{tabular}{lcc}
\toprule
\textbf{Counting mode} & \textbf{nRF ($\downarrow$)} & \textbf{Matches reference?} \\
\midrule
Canonical & 0.17 & No \\
Non-canonical & \textbf{0.00} & Yes \\
\bottomrule
\end{tabular}
\end{table}

\subsection{Overall ranking across all seven datasets}

Aggregating performance across all seven datasets, Peafowl achieved the best nRF on three datasets (7 Primates, E.~coli/Shigella 29, Plants), while phylonium led on four (Table~\ref{tab:summary_ranking}). Peafowl's average nRF across all datasets was 0.116 in canonical mode, compared to 0.119 for phylonium, indicating comparable overall accuracy.

\begin{table}[H]
\centering
\caption{Summary ranking of Peafowl across all seven datasets. Mean nRF ($\downarrow$ lower is better) is computed across datasets for each method. Peafowl non-canonical result used for Yersinia. Wilcoxon signed-rank test comparing Peafowl vs.\ each competitor across seven datasets.}
\label{tab:summary_ranking}
\begin{tabular}{lccc}
\toprule
\textbf{Method} & \textbf{Mean nRF ($\downarrow$)} & \textbf{Datasets best} & \textbf{$p$-value vs.\ Peafowl} \\
\midrule
\textbf{Peafowl} & \textbf{0.091} & \textbf{3} & --- \\
phylonium & 0.098 & 4 & 0.813 \\
andi & 0.138 & 0 & 0.047 \\
Skmer & 0.157 & 0 & 0.031 \\
FSWM & 0.197 & 0 & 0.016 \\
co-phylog & 0.289 & 0 & 0.016 \\
Multi-SpaM & 0.200 & 0 & 0.016 \\
Mash & 0.246 & 0 & 0.016 \\
FFP & 0.293 & 0 & 0.016 \\
CAFE-cvtree & 0.367 & 0 & 0.016 \\
\bottomrule
\end{tabular}
\end{table}

\subsection{Runtime and memory usage}

Table~\ref{tab:runtime} reports the runtime and peak memory consumption of Peafowl across all datasets. Peafowl's runtime scales with both genome size and number of taxa, primarily due to the k-mer generation step (Jellyfish) and the RAxML ML optimization.

\begin{table}[H]
\centering
\caption{Runtime (wall-clock) and peak memory usage of Peafowl on each benchmark dataset. Measurements obtained on a server with 64-core AMD EPYC 7763 and 512 GB RAM.}
\label{tab:runtime}
\begin{tabular}{lccc}
\toprule
\textbf{Dataset} & \textbf{Taxa} & \textbf{Runtime (min)} & \textbf{Peak memory (GB)} \\
\midrule
7 Primates & 7 & 2.3 & 0.8 \\
Drosophila & 14 & 47.6 & 12.4 \\
E.~coli/Shigella (29) & 29 & 18.9 & 4.2 \\
Labroidei & 25 & 8.4 & 2.1 \\
Plants & 14 & 124.5 & 38.7 \\
E.~coli/Shigella (27) & 27 & 16.7 & 3.9 \\
Yersinia & 8 & 3.1 & 1.2 \\
\bottomrule
\end{tabular}
\end{table}

\section{Discussion}

We have presented Peafowl, the first alignment-free phylogenetic tool that bridges k-mer-based data preparation with maximum likelihood tree estimation. By encoding k-mer presence/absence as a binary character matrix, Peafowl transforms the alignment-free paradigm from a purely distance-based framework into one compatible with the statistically rigorous ML inference machinery of RAxML \citep{stamatakis2014}. Across seven benchmark datasets, Peafowl achieved the lowest nRF distance on three datasets and was competitive on the remaining four, demonstrating that AF-ML inference can match or exceed the accuracy of established AF-distance methods.

The entropy-based k-mer selection strategy is a key practical contribution. By identifying the knee point in the cumulative entropy curve, Peafowl automatically determines an appropriate k-mer length without requiring user intervention or exhaustive grid search. This addresses a persistent challenge in AF methods, where k-mer length is often set heuristically or left as a user-tunable parameter \citep{zielezinski2019}. Our analysis confirms that the selected $k_{\text{entropy}}$ consistently yields nRF values at or near the optimum for each dataset.

The non-canonical counting mode provides a notable advantage for datasets containing genomic inversions. On the Yersinia dataset, canonical counting missed inversion events and produced a non-zero nRF, while non-canonical counting recovered the reference tree exactly. This capability is absent from most competing AF methods that either implicitly use canonical counting or do not expose the choice to users.

Several limitations warrant discussion. First, Peafowl currently requires assembled sequences or genomes and cannot directly process raw sequencing reads. Second, the binary encoding discards k-mer count information, potentially losing signal from copy-number variation. Third, the current implementation scales to approximately 30 taxa; larger datasets would require matrix dimensionality reduction or approximate ML methods. Fourth, performance degrades when comparing highly divergent species that share few k-mers, reducing the effective character set for ML inference.

Despite these limitations, our results establish a proof of concept that AF data are amenable to ML phylogenetic estimation. Future work should explore count-based encodings, support for unassembled reads via streaming k-mer analysis, scalability improvements for larger taxon sets, and the integration of model selection procedures analogous to those used in alignment-based ML phylogenetics \citep{huelsenbeck2001}.

\section{Methods}

\subsection{K-mer generation}

For each input genome, k-mers of odd lengths $k \in \{9, 11, 13, \ldots, 31\}$ were enumerated using Jellyfish version 2.3.0 \citep{marcais2011}. Both canonical and non-canonical counting modes were supported. In canonical mode, a k-mer and its reverse complement are treated as identical; in non-canonical mode, they are counted separately. Hash table sizes were set to $10^8$ for genomes $< 100$ Mb and $10^9$ otherwise.

\subsection{Binary matrix construction}

For each k-mer length, a binary matrix $\mathbf{M} \in \{0,1\}^{K \times N}$ was constructed, where $K$ is the total number of distinct k-mers observed across all $N$ taxa, and $M_{ij} = 1$ if k-mer $i$ is present in taxon $j$. K-mers present in all taxa or in only one taxon were removed as uninformative.

\subsection{Entropy-based k-mer length selection}

For each binary matrix at k-mer length $k$, the Shannon entropy per column was computed as $H_j(k) = -\sum_{i} p_i \log_2 p_i$, where $p_i$ is the fraction of k-mers present (or absent) in taxon $j$. The cumulative entropy $C(k) = \sum_{j=1}^{N} H_j(k)$ was plotted against $k$, and the knee point---identified via the Kneedle algorithm---was selected as $k_{\text{entropy}}$.

\subsection{Maximum likelihood tree estimation}

The binary matrix at $k_{\text{entropy}}$ was provided as a PHYLIP-format alignment to RAxML version 8.2.12 \citep{stamatakis2014} using the BINGAMMA model (binary data with gamma-distributed rate heterogeneity). One hundred bootstrap replicates were performed, and the best-scoring ML tree was selected.

\subsection{Benchmarking}

Nine AF tools were benchmarked: FFP \citep{sims2009}, co-phylog \citep{yi2013}, Mash \citep{ondov2016}, Skmer \citep{sarmashghi2019}, FSWM \citep{leimeister2017}, andi \citep{haubold2015}, phylonium \citep{kloetzl2020}, Multi-SpaM \citep{dencker2020}, and CAFE-cvtree \citep{lu2017}. Distance matrices from these tools were converted to neighbor-joining trees using MEGA-X \citep{kumar2018}. Robinson--Foulds distances were computed with PHYLIP \citep{felsenstein2005} and normalized by the maximum possible RF distance: $\text{nRF} = \text{RF} / (2n - 6)$, where $n$ is the number of taxa.

\subsection{Statistical analysis}

The Wilcoxon signed-rank test was used to compare Peafowl's nRF values against each competing method across the seven datasets, with significance set at $\alpha = 0.05$. No correction for multiple comparisons was applied given the exploratory nature of the analysis.

\section{Conclusion}

Peafowl demonstrates that alignment-free phylogenetic data, when encoded as binary k-mer presence/absence matrices, can be effectively analyzed using maximum likelihood inference. The entropy-based k-mer selection provides a principled, automatic approach to a parameter that has historically required manual tuning. Across seven diverse benchmark datasets, Peafowl achieved competitive or superior accuracy compared to nine established AF methods, achieving perfect reference-tree recovery on two datasets. These results open a new direction for alignment-free phylogenomics by bringing the statistical power of likelihood-based methods to bear on k-mer-derived data.

\bibliographystyle{plainnat}
\bibliography{references}

\end{document}
